\preludeandfugue{Pipe, glass, bottle}
\contributor{Sadie Kaye}

\begin{flushright}
\textit{Form strives to liberate itself from matter\ldots}\\
\'Etienne Gilson, \textit{Painting and reality},\\
Princeton University Press, 1957.
\end{flushright}

\begin{flushright}
\textit{Pipe, verre, bouteille de Vieux Marc.}\\
Pablo Picasso, 1914.
\end{flushright}

% abstract vs concrete
% material vs virtual
% transformed, repurposed, appropriated

\prelude

\begin{verse}

O pluck out the Spanish songs\\
on this sketchiest of false guitars,\\
and sing out its transformation!\\!

Sing out the thunderous vocation\\
of a torn and ragg\`ed sheet,\\
a single page of frozen newsprint,\\!

and pipe your dreams through the work\\
with the scantest splash of v\'eritable vieux Marc!\\
Sant\'e!  Bottoms up!\\!

\end{verse}

\fugue

\clearpage

\begin{verse}

Matter seeks shape to appropriate\\
and mould itself to living taste:\\
so the artist acts---\\!

and art is matter's form:\\
this material and form are proximate\\
yet form will strive to liberate\\!

itself from such worldly stuff;\\
though, without matter's grist,\\
it cannot truly subsist.\\!

In this way art's substance may subsume\\
anything it could desire,\\
because art can,\\!

and should the artist expropriate\\
primary matter from a futurist rag\\
pasted down on board\\!

this cheap newsprint will then become\\
the boldest of icons, preserved forever \\
to grace a different world.\\!

\end{verse}
